\documentclass[11pt]{article}

\usepackage[margin=1in]{geometry}
\usepackage{booktabs}
\usepackage{amsmath}
\usepackage{amssymb}
\usepackage{hyperref}
\usepackage{cite}
\usepackage{authblk}

\title{Clustered Synapses Develop in Distinct Dendritic Domains in the Visual Cortex Before Eye Opening}

\author[1]{Anonymous Author}
\affil[1]{Anonymous Institution}

\date{}

\begin{document}

\maketitle

\begin{abstract}
The organization of synapses along the dendrites of cortical neurons shapes how those neurons integrate their inputs. It has long been proposed that clustering of co-active synaptic inputs on dendrites enhances the computational power of single neurons, but when and how such clustering arises during development remains incompletely understood. Here we use in utero electroporation of GCaMP6s and DsRed into layer 2/3 pyramidal neurons of mouse visual cortex, combined with two-photon imaging of individual synapses ex vivo during the second postnatal week ($P8$--$P13$), to map the functional organization of synaptic inputs along dendrites before eye opening. We find that the frequency of synaptic transmission per synapse rises from $P8$--$10$ to $P12$--$13$, with spine synapses systematically more active than shaft synapses at both ages ($P8$--$10$: $0.55 \pm 0.09$\,min$^{-1}$ vs.\ $0.16 \pm 0.03$\,min$^{-1}$; $P12$--$13$: $0.67 \pm 0.07$\,min$^{-1}$ vs.\ $0.23 \pm 0.03$\,min$^{-1}$; all $p < 10^{-9}$, Mann-Whitney U). The majority of synaptic inputs occur in barrages ($59\%$ at $P8$--$10$, $61\%$ at $P12$--$13$). Most active synapses are located on spines ($76\%$ of high-activity synapses), and synapses within $10\,\mu\mathrm{m}$ of a high-activity synapse exhibit higher mean transmission frequency than those farther away ($n = 181$ vs.\ $173$ synapses, independent-samples $t$-test, $2$-sided). We define dendritic domains as segments surrounding high-activity synapses, and show that both the number of synapses per domain (from $3.2 \pm 1.7$ to $6.0 \pm 3.2$) and the density of domains along dendrites (from $1.4 \pm 0.7$ to $2.6 \pm 1.0$ per $100\,\mu\mathrm{m}$) approximately double during the second postnatal week, with domains covering $27.2 \pm 16.2\%$ of dendritic length at $P8$--$10$ and $61.3 \pm 29.4\%$ at $P12$--$13$. The percentage of synapses within domains rises from $46.2 \pm 7.4\%$ to $72 \pm 23.2\%$ ($n = 6$ dendrites per age, $p = 0.026$, $t$-test). Local co-activity is higher within domains than outside ($P8$--$10$: $0.08 \pm 0.02$ vs.\ $0.02 \pm 0.01$; $P12$--$13$: $0.06 \pm 0.01$ vs.\ $0.04 \pm 0.01$), and Mann-Kendall scores of monotonic change in transmission frequency are higher inside domains. Together these data establish that functional synaptic clustering emerges in distinct dendritic domains in visual cortex before eye opening.
\end{abstract}

\section{Introduction}

A major contributor to the computational capacity of cortical neurons is the spatial arrangement of synaptic inputs along their dendrites \cite{hausser2001diversity,poirazi2003pyramidal,larkum2004active}. Clustered, co-active synapses on a single dendritic branch can recruit supralinear dendritic nonlinearities and local plasticity rules that operate at finer spatial scales than somatic activity \cite{losonczy2008compartment,golding2002dendritic,govindarajan2011memory,magee1997synaptic}. Evidence from adult hippocampus and cortex indicates that synapses with similar functional properties tend to be clustered along short stretches of dendrite \cite{iacaruso2017synaptic,scholl2017local,wilson2016division,kerlin2019functional,fu2012repetitive}, suggesting that clustering may be an organizational principle of cortical circuits \cite{yuste2011dendritic}.

When clustering arises during development is less clear. In the developing hippocampus, spontaneous activity drives local synaptic plasticity in vivo \cite{winnubst2015} and functional synaptic clusters can already be detected \cite{kleindienst2011clustered}. In the visual cortex, spontaneous activity in the week preceding eye opening is organized into barrages of correlated network events \cite{garaschuk2000calcium,khazipov2004early,siegel2012peripheral,allene2008sequential,cline2008spontaneous,katz1996synaptic}, and this activity is thought to shape circuit wiring before pattern vision. Whether the spatial arrangement of synaptic inputs on individual layer 2/3 cortical neurons becomes clustered during this pre-vision period, and whether such clustering is organized into discrete dendritic domains, has not been mapped.

Here we directly address these questions by imaging large numbers of spontaneously active individual synapses at single-synapse resolution along dendrites of layer 2/3 pyramidal neurons of mouse visual cortex across the second postnatal week ($P8$--$P13$), before eye opening. We show that functional synaptic clusters organize into spatially distinct dendritic domains, that both the number and density of domains approximately double during the second postnatal week, and that co-activity of synapses is strongly enhanced within domains.

\section{Methods}

\paragraph{In utero electroporation.} Pyramidal neurons in layer 2/3 of the visual cortex were co-transfected with GCaMP6s \cite{chen2013gcamp6} ($2$\,mg/ml) and DsRed ($0.5$--$2$\,mg/ml) at E16.5 using in utero electroporation following Winnubst et al.\ \cite{winnubst2015,saito2001invivo}. Pregnant mice were anesthetized with $2\%$ isoflurane (reduced to $0.7$--$1\%$ after surgery) and a small ($1.5$--$2$\,cm) incision was made in the abdominal wall. The uterine horns were removed from the abdomen, and $1\,\mu\mathrm{l}$ of DNA was injected into the lateral ventricle of each embryo using a sharp glass pipette. Voltage pulses (five square-wave pulses, $30$\,V, $50$\,ms duration, $950$\,ms interval; custom-built electroporator) were delivered across the brain with tweezer electrodes covered in conductive gel.

\paragraph{Electrophysiology.} To visualize synaptic inputs, action potentials were blocked with QX-314 in the intracellular solution ($120$\,mM CsMeSO$_3$, $8$\,mM NaCl, $15$\,mM CsCl$_2$, $10$\,mM TEA-Cl, $10$\,mM HEPES, $5$\,mM QX-314 bromide, $4$\,mM MgATP, and $0.3$\,mM Na-GTP) \cite{takahashi2012,winnubst2015}. Currents were recorded in voltage-clamp mode at $10$\,kHz and filtered at $3$\,kHz (Multiclamp 700B; Molecular Devices).

\paragraph{Two-photon imaging of synapses.} Dendrites of GCaMP6s/DsRed-labelled layer 2/3 pyramidal neurons were imaged at $5$--$15$\,Hz with a pixel size of $0.32\,\mu\mathrm{m}$ in single planes or small stacks with planes spaced $1.5$--$2\,\mu\mathrm{m}$. Recording durations averaged $29 \pm 8$\,min at $P8$--$10$ and $26 \pm 6$\,min at $P12$--$13$ (mean $\pm$ STD).

\paragraph{Definition of dendritic domains.} Synapses whose transmission frequency was in the top $20\%$ within their age group were classified as high-activity synapses. Dendritic domains were defined as segments of dendrite surrounding high-activity synapses. Synapse clusters were defined as groups of neighbouring synapses with similar activity patterns. Randomised distributions of synapse positions were constructed by assigning ``empty'' positions between the actual functional-synapse positions such that the entire imaged dendrite was covered with real or empty positions at the typical minimal inter-synapse distance ($4 \pm 1\,\mu\mathrm{m}$). Inactive structural spines (potentially presynaptically silent) made up $25 \pm 20\%$ of spines at $P8$--$10$ and $9 \pm 9\%$ at $P12$--$13$ (mean $\pm$ STD) \cite{voronin2004silent}.

\paragraph{Statistics.} Per-synapse comparisons used Mann-Whitney U tests (2-sided), independent-samples Student's $t$-tests (2-sided), Wilcoxon signed-rank tests, Mann-Kendall trend tests \cite{mann2019kendall} and linear regression; $95\%$ confidence intervals on regressions are shown where relevant. Statistical parameters from Fig.~3 were derived from linear regressions on the data from $n = 12$ experiments.

\section{Results}

\subsection{Synaptic transmission becomes more frequent by the end of the second postnatal week}
Comparison of $n = 6$ dendrites at $P8$--$10$ with $n = 6$ dendrites at $P12$--$13$ showed that the number of synaptic transmission events per dendrite was higher at $P12$--$13$. Linear regression across all $n = 12$ experiments confirmed that the per-synapse frequency of synaptic transmission rose with age. At both ages, spine synapses were more active than shaft synapses (Table~\ref{tab:spineshaft}; all $p < 10^{-9}$, Mann-Whitney U; $P8$--$10$: $n = 36$ spine / $113$ shaft synapses; $P12$--$13$: $n = 83$ spine / $122$ shaft synapses). The distribution of transmission frequencies was long-tailed, with most synapses receiving inputs at low frequencies and a minority being very active. A fraction of structural spines remained inactive during recordings ($25 \pm 20\%$ at $P8$--$10$, $9 \pm 9\%$ at $P12$--$13$), consistent with a presynaptically silent population \cite{voronin2004silent}.

\begin{table}[h]
\centering
\caption{Mean $\pm$ SEM synaptic transmission frequency (min$^{-1}$) at spine vs.\ shaft synapses.}
\label{tab:spineshaft}
\begin{tabular}{lcc}
\toprule
Age & Spine synapses & Shaft synapses \\
\midrule
$P8$--$10$ & $0.55 \pm 0.09$ & $0.16 \pm 0.03$ \\
$P12$--$13$ & $0.67 \pm 0.07$ & $0.23 \pm 0.03$ \\
\bottomrule
\end{tabular}
\end{table}

\subsection{Synaptic transmission is organized in barrages}
Most synaptic inputs occurred in barrages: $59\%$ of events at $P8$--$10$ and $61\%$ at $P12$--$13$ (Fig.~4A), consistent with the correlated network events that dominate spontaneous activity in visual cortex during this developmental window \cite{siegel2012peripheral,allene2008sequential,khazipov2004early,garaschuk2000calcium}.

\subsection{High-activity synapses define dendritic domains}
We classified the top-$20\%$ most active synapses within each age group as high-activity synapses (Fig.~5A). The majority ($76\%$) of high-activity synapses were on spines, consistent with spine synapses being the more active population (Table~\ref{tab:spineshaft}). Mean transmission frequency at synapses within $10\,\mu\mathrm{m}$ of a high-activity synapse was significantly higher than at synapses farther away (Student's $t$-test for independent samples, 2-sided; $n = 181$ synapses within $10\,\mu\mathrm{m}$ and $n = 173$ synapses beyond $10\,\mu\mathrm{m}$).

\subsection{Domain properties change during the second postnatal week}
Mean domain extent grew modestly during the second postnatal week (from $18.7 \pm 6.1\,\mu\mathrm{m}$ at $P8$--$10$ to $22.5 \pm 6.4\,\mu\mathrm{m}$ at $P12$--$13$). The number of synapses per domain approximately doubled (from $3.2 \pm 1.7$ to $6.0 \pm 3.2$), as did the density of domains along dendrites (from $1.4 \pm 0.7$ to $2.6 \pm 1.0$ per $100\,\mu\mathrm{m}$). The fraction of dendritic length covered by domains increased significantly across this interval, from $27.2 \pm 16.2\%$ at $P8$--$10$ to $61.3 \pm 29.4\%$ at $P12$--$13$ (mean $\pm$ STD). The percentage of synapses within domains rose from $46.2 \pm 7.4\%$ at $P8$--$10$ ($n = 6$ dendrites) to $72 \pm 23.2\%$ at $P12$--$13$ ($n = 6$ dendrites; mean $\pm$ STD; $p = 0.026$, $t$-test). Table~\ref{tab:domains} summarises these domain statistics.

\begin{table}[h]
\centering
\caption{Domain properties at $P8$--$10$ vs.\ $P12$--$13$ (mean $\pm$ STD where indicated).}
\label{tab:domains}
\begin{tabular}{lcc}
\toprule
Metric & $P8$--$10$ & $P12$--$13$ \\
\midrule
Domain extent ($\mu$m) & $18.7 \pm 6.1$ & $22.5 \pm 6.4$ \\
Synapses per domain & $3.2 \pm 1.7$ & $6.0 \pm 3.2$ \\
Domain density (per $100\,\mu$m) & $1.4 \pm 0.7$ & $2.6 \pm 1.0$ \\
Dendritic length covered by domains (\%) & $27.2 \pm 16.2$ & $61.3 \pm 29.4$ \\
Synapses within domains (\%) & $46.2 \pm 7.4$ & $72 \pm 23.2$ \\
\bottomrule
\end{tabular}
\end{table}

\subsection{Synapses within domains co-activate more strongly than synapses outside}
Local co-activity of synapses was higher within domains than outside: at $P8$--$10$, within $0.08 \pm 0.02$ vs.\ outside $0.02 \pm 0.01$; at $P12$--$13$, within $0.06 \pm 0.01$ vs.\ outside $0.04 \pm 0.01$ (mean $\pm$ SEM; Fig.~6F). Pooling across ages, the local co-activity of synapses inside domains was higher than that of outside-domain synapses (Mann-Whitney U, 2-sided; $n = 213$ in-domain vs.\ $173$ outside-domain synapses). For each experiment, the co-activity of synapses within domains was higher than their co-activity with synapses in other domains (Wilcoxon signed-rank test).

Mann-Kendall scores, quantifying monotonic changes in per-synapse transmission frequency, were higher for domain synapses than for synapses outside domains (Student's $t$-test for independent samples, 2-sided; $n = 213$ in-domain vs.\ $173$ outside-domain synapses from all ages).

\section{Discussion}

By imaging hundreds of individual synapses along the dendrites of layer 2/3 pyramidal neurons in mouse visual cortex before eye opening, we show that functional synaptic clustering emerges in spatially distinct dendritic domains during the second postnatal week. Three observations support this conclusion. First, synaptic transmission frequency increases from $P8$--$10$ to $P12$--$13$, with spine synapses systematically more active than shaft synapses ($0.55 \pm 0.09$ vs.\ $0.16 \pm 0.03$\,min$^{-1}$ at $P8$--$10$; $0.67 \pm 0.07$ vs.\ $0.23 \pm 0.03$\,min$^{-1}$ at $P12$--$13$; all $p < 10^{-9}$). Second, most synaptic inputs occur in barrages, and $76\%$ of the high-activity synapses lie on spines. Synapses within $10\,\mu\mathrm{m}$ of a high-activity synapse are significantly more active than synapses further away, suggesting a local influence of high-activity synapses on neighbouring synapses. Third, defining dendritic domains as segments surrounding high-activity synapses reveals that both the number and density of domains approximately double between $P8$--$10$ and $P12$--$13$, that domains cover $27.2 \pm 16.2\%$ of dendritic length at $P8$--$10$ and $61.3 \pm 29.4\%$ at $P12$--$13$, and that the fraction of synapses within domains rises from $46.2 \pm 7.4\%$ to $72 \pm 23.2\%$ ($p = 0.026$, $t$-test).

Several lines of evidence support the biological relevance of this domain-centred organization. Co-activity is higher inside than outside domains in both younger ($0.08 \pm 0.02$ vs.\ $0.02 \pm 0.01$) and older animals ($0.06 \pm 0.01$ vs.\ $0.04 \pm 0.01$), and remains elevated after controlling for comparisons between domains within the same dendrite (Wilcoxon signed-rank test). Furthermore, the Mann-Kendall trend score of transmission-frequency change is higher for synapses inside domains, indicating that domains preferentially host synapses whose activity is monotonically changing --- a hallmark of plasticity.

Our results extend prior observations of activity-dependent clustering of functional synaptic inputs in the hippocampus \cite{kleindienst2011clustered,winnubst2015} to the developing mouse visual cortex before eye opening, and complement adult studies showing functional synaptic clustering along dendrites \cite{iacaruso2017synaptic,scholl2017local,wilson2016division,kerlin2019functional,fu2012repetitive}. They suggest that clustered-input architectures proposed to be essential for supralinear dendritic computation \cite{poirazi2003pyramidal,larkum2004active,golding2002dendritic,govindarajan2011memory} begin to emerge during the week preceding the onset of pattern vision, driven presumably by spontaneously generated barrages of activity \cite{katz1996synaptic,siegel2012peripheral,allene2008sequential,khazipov2004early}. Limitations include the acute nature of the slice preparation, the relatively small number of dendrites ($n = 6$ per age), and our focus on mouse visual cortex; in vivo, multi-day longitudinal imaging would be needed to test whether the same dendritic domains persist across development. Beyond these caveats, our finding that clustered synapses develop in distinct dendritic domains in the visual cortex before eye opening provides a principled anatomical and functional substrate for dendritic computation in the developing neocortex.

\bibliographystyle{unsrt}
\bibliography{references}

\end{document}
